\ifdefined\gkver\else\def\gkver{student}\fi
\documentclass[\gkver]{gaokaozhenti}
\begin{document}

\chapter{2010年海南}

\begin{enumerate}
\item
%% number: 1
%% paperName: 2010年海南高考第 1 题 3 分
%% typeId: 1
%% type: 单选题
%% chapter: 磁场
%% point: 测定直流电动机的效率
%% method: 无
%% score: 3
%% degree: 650
%% duplicateId: 0
%% body:
1873年奥地利维也纳世博会上，比利时出生的法国工程师格拉姆在布展中偶然接错了导线，把另一直流发电机发出的电接到了自己送展的直流发电机的电流输出端。由此而观察到的现象导致了他的一项重要发明，从而突破了人类在电能利用方中的一个瓶颈。此项发明是\xzanswer{B}

\fourchoices[answer=B]
{新型直流发电机}
{直流电动机}
{交流电动机}
{交流发电机}

%% answer: B
%% memo:
\memoanswer{
直流发电机发电时接另一直流发电机，则另一直流发电机实际成了直流电动机，B 正确。
}

\item
%% number: 2
%% paperName: 2010年海南高考第 2 题 3 分
%% typeId: 1
%% type: 单选题
%% chapter: 电磁感应
%% point: 楞次定律推广
%% method: 无
%% score: 3
%% degree: 650
%% duplicateId: 0
%% body:
一金属圆环水平固定放置。现将一竖直的条形磁铁，在圆环上方沿圆环轴线从静止开始释放，在条形磁铁穿过圆环的过程中，条形磁铁与圆环\xzanswer{D}

\fourchoices[answer=D]
{始终相互吸引}
{始终相互排斥}
{先相互吸引，后相互排斥}
{先相互排斥，后相互吸引}

%% answer: D
%% memo:
\memoanswer{
由楞次定律可知，当条形磁铁靠近圆环时，感应电流阻碍其靠近，是排斥力；当磁铁穿过圆环远离圆环时，感应电流阻碍其远离，是吸引力，D 正确。
}

\item
%% number: 3
%% paperName: 2010年海南高考第 3 题 3 分
%% typeId: 1
%% type: 单选题
%% chapter: 直线运动
%% point: 匀速直线运动
%% method: 无
%% score: 3
%% degree: 650
%% duplicateId: 0
%% body:
下列说法正确的是\xzanswer{D}

\fourchoices[answer=D]
{若物体运动速率始终不变，则物体所受合力一定为零}
{若物体的加速度均匀增加，则物体做匀加速直线运动}
{若物体所受合力与其速度方向相反，则物体做匀减速直线运动}
{若物体在任意的相等时间间隔内位移相等，则物体做匀速直线运动}

%% answer: D
%% memo:
\memoanswer{
物体运动速率不变但方向可能变化，因此合力不一定为零，A 错；物体的加速度均匀增加，即加速度在变化，是非匀加速直线运动，B 错；物体所受合力与其速度方向相反，只能判断其做减速运动，但加速度大小不可确定，C 错；若物体在任意的相等时间间隔内位移相等，则物体做匀速直线运动，D 对。
}

\item
%% number: 4
%% paperName: 2010年海南高考第 4 题 3 分
%% typeId: 1
%% type: 单选题
%% chapter: 电场
%% point: 电场叠加
%% method: 无
%% score: 3
%% degree: 450
%% duplicateId: 0
%% body:
如右图，M、N和P是以MN为直径的半圆弧上的三点，O点为半圆弧的圆心，$\angle$MOP$=60^{\circ}$。电荷量相等、符号相反的两个点电荷分别置于M、N两点，这时O点电场强度的大小为$E_{1}$；若将N点处的点电荷移至P点，则O点的场强大小变为$E_{2}$，$E_{1}$与$E_{2}$之比为\xzanswer{B}

\onepicture[align=right, w=4.5cm]{figs/04.png}

\fourchoices[answer=B]
{1$:$2}
{2$:$1}
{2$:\sqrt{3}$}
{4$:\sqrt{3}$}

%% answer: B
%% memo:
\memoanswer{
依题意，每个点电荷在 O 点产生的场强为 $\frac{E_{1}}{2}$，则当 N 点处的点电荷移至 P 点时，O 点场强如图所示，合场强大小为 $E_{2}=\frac{E_{1}}{2}$，则 $\frac{E_{1}}{E_{2}}=\frac{2}{1}$，B 正确。

\onepicture[w=4.0cm]{figs/04_an.png}
}

\item
%% number: 5
%% paperName: 2010年海南高考第 5 题 3 分
%% typeId: 1
%% type: 单选题
%% chapter: 运动和力
%% point: 牛顿第二定律
%% method: 无
%% score: 3
%% degree: 450
%% duplicateId: 0
%% body:
如右图，水平地面上有一楔形物块 a，其斜面上有一小物块 b，b 与平行于斜面的细绳的一端相连，细绳的另一端固定在斜面上。a 与 b 之间光滑，a 和 b 以共同速度在地面轨道的光滑段向左运动。当它们刚运行至轨道的粗糙段时\xzanswer{C}

\onepicture[align=right, w=4.8cm]{figs/05.png}

\fourchoices[answer=C]
{绳的张力减小，b 对 a 的正压力减小}
{绳的张力增加，斜面对 b 的支持力增加}
{绳的张力减小，地面对 a 的支持力增加}
{绳的张力增加．地面对 a 的支持力减小}

%% answer: C
%% memo:
\memoanswer{
在光滑段运动时，系统及物块 b 处于平衡状态，因此有 $T\cos\theta-N\sin\theta=0$，$T\sin\theta+N\cos\theta-mg=0$；当它们刚运行至轨道的粗糙段时，系统有水平向右的加速度，此时有两种可能，一是物块 b 仍相对静止，竖直方向加速度为零，则 $T\sin\theta+N\cos\theta-mg=0$ 仍成立，但 $T\cos\theta-N\sin\theta=ma<0$，因此绳的张力 $T$ 将减小，而地面对 $a$ 的支持力不变；二是物块 b 相对于 a 向上滑动，具有向上的加速度，是超重，因此绳的张力减小，地面对 a 的支持力增大，C 正确。

\onepicture[w=3.2cm]{figs/05_an.png}
}

\item
%% number: 6
%% paperName: 2010年海南高考第 6 题 3 分
%% typeId: 1
%% type: 单选题
%% chapter: 运动和力
%% point: 牛顿定律的应用
%% method: 无
%% score: 3
%% degree: 450
%% duplicateId: 0
%% body:
在水平的足够长的固定木板上，一小物块以某一初速度开始滑动，经一段时间 $t$ 后停止。现将该木板改置成倾角为 $45^{\circ}$ 的斜面，让小物块以相同的初速度沿木板上滑。若小物块与木板之间的动摩擦因数为 $\mu$，则小物块上滑到最高位置所需时间与 $t$ 之比为\xzanswer{A}

\fourchoices[answer=A]
{$\frac{\sqrt{2}}{1+\mu}$}
{$\frac{\mu}{1+\sqrt{2}\mu}$}
{$\frac{\mu}{\sqrt{2}+\mu}$}
{$\frac{1+\mu}{\sqrt{2}\mu}$}

%% answer: A
%% memo:
\memoanswer{
木板水平时，小物块的加速度 $a_{1}=\mu g$，设滑行初速度为 $v_{0}$，则滑行时间为 $t=\frac{v_{0}}{\mu g}$；木板改成后，小物块上滑的加速度 $a_{2}=\frac{mg\sin 45^{\circ}+\mu mg\cos 45^{\circ}}{m}=\frac{(1+\mu)\sqrt{2}g}{2}$，滑行时间 $t'=\frac{v_{0}}{a_{2}}=\frac{v_{0}}{(1+\mu)\sqrt{2}g}$，因此 $\frac{t'}{t}=\frac{a_{1}}{a_{2}}=\frac{\sqrt{2}\mu}{1+\mu}$，A 项正确。
}

\item
%% number: 7
%% paperName: 2010年海南高考第 7 题 4 分
%% typeId: 2
%% type: 多选题
%% chapter: 电磁感应
%% point: 自感与互感，涡流
%% method: 无
%% score: 4
%% degree: 650
%% duplicateId: 0
%% body:
下列说法正确的是\xzanswer{AC}

\fourchoices[answer=AC]
{当线圈中电流不变时，线圈中没有自感电动势}
{当线圈中电流反向时，线圈中自感电动势的方向与线圈中原电流的方向相反}
{当线圈中电流增大时，线圈中自感电动势的方向与线圈中电流的方向相反}
{当线圈中电流减小时，线圈中自感电动势的方向与线圈中电流的方向相反}

%% answer: AC
%% memo:
\memoanswer{
由法拉第电磁感应定律可知，当线圈中电流不变时，不产生自感电动势，A 对；

当线圈中电流反向时，相当于电流减小，线圈中自感电动势的方向与线圈中原电流的方向相同，B 错；

当线圈中电流增大时，自感电动势阻碍电流的增大，线圈中自感电动势的方向与线圈中电流的方向相反，C 对；

当线圈中电流减小时，自感电动势阻碍电流的减小，线圈中自感电动势的方向与线圈中电流的方向相同，D 错。
}

\item
%% number: 8
%% paperName: 2010年海南高考第 8 题 4 分
%% typeId: 2
%% type: 多选题
%% chapter: 运动和力
%% point: 超重失重
%% method: 无
%% score: 4
%% degree: 650
%% duplicateId: 0
%% body:
如右图，木箱内有一竖直放置的弹簧，弹簧上方有一物块，木箱静止时弹簧处于压缩状态且物块压在箱顶上。若在某一段时间内，物块对箱顶刚好无压力。则在此段时间内，木箱的运动状态可能为\xzanswer{BD}

\onepicture[align=right, w=2.4cm]{figs/08.png}

\fourchoices[answer=BD]
{加速下降}
{加速上升}
{减速上升}
{减速下降}

%% answer: BD
%% memo:
\memoanswer{
木箱静止时物块对箱顶有压力，则物块受到顶向下的压力，当物块对箱顶刚好无压力时，表明系统有向上的加速度，是超重，BD 正确。
}

\item
%% number: 9
%% paperName: 2010年海南高考第 9 题 4 分
%% typeId: 2
%% type: 多选题
%% chapter: 交流电
%% point: 变压器
%% method: 无
%% score: 4
%% degree: 650
%% duplicateId: 0
%% body:
如右图，一理想变压器原副线圈匝数之比为 $4:1$，原线圈两端接入一正弦交流电源；副线圈电路中 $R$ 为负载电阻，交流电压表和交流电流表都是理想电表。下列结论正确的是\xzanswer{AD}

\onepicture[align=right, w=4.5cm]{figs/09.png}

\fourchoices[answer=AD]
{若电压表读数为 $6\UV$，则输入电压的最大值为 $24\sqrt{2}\UV$}
{若输入电压不变，副线圈匝数增加到原来的 2 倍，则电流表的读数减小到原来的一半}
{若输入电压不变，负载电阻的阻值增加到原来的 2 倍，则输入功率也增加到原来的 2 倍}
{若保持负载电阻的阻值不变．输入电压增加到原来的 2 倍，则输出功率增加到原来的 4 倍}

%% answer: AD
%% memo:
\memoanswer{
若电压表读数为 $6\UV$，则输入电压为 $U_{1}=\frac{4}{1}\times6\UV=24\UV$ 是有效值，因此其最大值为 $24\sqrt{2}\UV$，A 正确；

若输入电压不变，副线圈匝数增加到原来的 2 倍，则输出电压也增加到原来的 2 倍，电流表示数应增加到原来的 2 倍，B 错；

若输入电压不变，负载电阻的阻值增加到原来的 2 倍，则输出电流减小到原来的一半，输入功率等于输出功率即 $P=IU$ 也减小到原来的一半，C 错；

若保持负载电阻的阻值不变，输入电压增加到原来的 2 倍，输出电压增大到原来的 2 倍，则由 $P=\frac{U^{2}}{R}$ 可知输出功率增加到原来的 4 倍，D 正确。
}

\item
%% number: 10
%% paperName: 2010年海南高考第 10 题 4 分
%% typeId: 2
%% type: 多选题
%% chapter: 曲线运动
%% point: 天体运动
%% method: 无
%% score: 4
%% degree: 450
%% duplicateId: 0
%% body:
火星直径约为地球的一半，质量约为地球的十分之一，它绕太阳公转的轨道半径约为地球公转半径的1.5倍。根据以上数据，以下说法正确的是\xzanswer{AB}

\fourchoices[answer=AB]
{火星表面重力加速度的数值比地球表面小}
{火星公转的周期比地球的长}
{火星公转的线速度比地球的大}
{火星公转的向心加速度比地球的大}

%% answer: AB
%% memo:
\memoanswer{
由 $G\frac{Mm}{R^{2}}=mg$ 得 $g=G\frac{M}{R^{2}}$，计算得火星表面的重力加速度约为地球表面的 $\frac{2}{5}$，A 正确；由 $G\frac{Mm}{r^{2}}=m(\frac{2\pi}{T})^{2}R$ 得 $T=2\pi\sqrt{\frac{r^{3}}{GM}}$，公转轨道半径大的周期长，B 对；周期长的线速度小（或由 $v=\sqrt{\frac{GM}{r}}$ 判断轨道半径大的线速度小），C 错；公转向心加速度 $a=G\frac{M}{r^{2}}$，D 错。
}

\item
%% number: 11
%% paperName: 2010年海南高考第 11 题 4 分
%% typeId: 3
%% type: 填空题
%% chapter: 电场
%% point: 静电
%% method: 无
%% score: 4
%% degree: 650
%% duplicateId: 0
%% body:
利用静电除尘器可以消除空气中的粉尘。静电除尘器由金属管A和悬在管中的金属丝B组成，A和B分别接到高压电源正极和负极，其装置示意图如右图所示。A、B之间有很强的电场，距B越近，场强\tkanswer[1cm]{越大}（填“越大”或“越小”）。B附近的气体分子被电离成为电子和正离子，粉尘吸附电子后被吸附到\tkanswer[1cm]{A}（填“A”或“B”）上，最后在重力作用下落入下面的漏斗中。

\onepicture[align=right, w=2.8cm]{figs/11.png}

%% answer: 
\jdanswer{
越大；A
}
%% memo:
\memoanswer{
电极截面如图所示，由电场线可判断越靠近 B 场强越大；粉尘吸附电子后带负电，因此向正极 A 运动。

\onepicture[w=3.0cm]{figs/11_an.png}
}

\item
%% number: 12
%% paperName: 2010年海南高考第 12 题 4 分
%% typeId: 3
%% type: 填空题
%% chapter: 运动和力
%% point: 变加速直线运动
%% method: 无
%% score: 4
%% degree: 250
%% duplicateId: 0
%% body:
雨滴下落时所受到的空气阻力与雨滴的速度有关，雨滴速度越大，它受到的空气阻力越大；此外，当雨滴速度一定时，雨滴下落时所受到的空气阻力还与雨滴半径的 $\alpha$ 次方成正比（$1\leqslant\alpha\leqslant2$）。假设一个大雨滴和一个小雨滴从同一云层同时下落，最终它们都\tkanswer[1.2cm]{匀速}（填“加速”、“减速”或“匀速”）下落。\tkanswer[1cm]{大}（填“大”或“小”）雨滴先落到地面；接近地面时，\tkanswer[1cm]{小}（填“大”或“小”）雨滴的速度较小。

%% answer: 
\jdanswer{
匀速；大；小
}
%% memo:
\memoanswer{
由于雨滴受到的空气阻力与速度有关，速度越大阻力越大，因此最终当阻力增大到与重力平衡时都做匀速运动；设雨滴半径为 $r$，则当雨滴匀速下落时受到的空气阻力 $f\propto r^{\alpha}$，而重力 $mg=\rho\frac{4}{3}\pi r^{3}$，由于 $1\leqslant\alpha\leqslant2$，因此半径大的匀速运动的速度大，先落地且落地速度大，小雨滴落地速度小。
}

\item
%% number: 13
%% paperName: 2010年海南高考第 13 题 6 分
%% typeId: 4
%% type: 实验题
%% chapter: 电路
%% point: 电路实验
%% method: 无
%% score: 6
%% degree: 650
%% duplicateId: 0
%% body:
图\ref{2010海南13a}为测量电压表V内阻 $r$ 的电路原理图。图中两个固定电阻的阻值均为 $R$，$\mathrm{S}_{1}$、$\mathrm{S}_{2}$是开关，$E$是电源（内阻可忽略）。

\twopicture[labelA=2010海南13a, labelB=2010海南13b, numA=1, numB=2, numstyle=arabic, wA=4.2cm, wB=4.8cm]
{figs/13a.png}
{figs/13b.png}

\begin{enumerate}
	\item
	按电路原理图将图\ref{2010海南13b}中的实物图连线；
	\item
	开关 $\mathrm{S}_{1}$ 保持断开，合上开关 $\mathrm{S}_{2}$，此时电压表的读数为 $U_{1}$；再合上开关 $\mathrm{S}_{1}$，电压表的读数变为 $U_{2}$，电压表内阻 $r=$\tkanswer[2cm]{$\frac{U_{1}}{U_{2}-2U_{1}}R$}（用 $U_{1}$、$U_{2}$ 和 $R$ 表示）。
\end{enumerate}

%% answer: 
\jdanswer{
\begin{enumerate}
	\item
	实物连线如图。

	\onepicture[w=6.5cm]{figs/13_an.png}

	\item
	$\frac{U_{1}}{U_{2}-2U_{1}}R$
\end{enumerate}
}
%% memo:
\memoanswer{
开关 $\mathrm{S}_{1}$ 保持断开时，$\frac{E}{U_{1}}=\frac{R+\frac{Rr}{R+r}}{\frac{Rr}{R+r}}$，合上开关 $\mathrm{S}_{1}$，则 $E=U_{2}$，解得 $r=\frac{U_{1}}{U_{2}-2U_{1}}R$。
}

\item
%% number: 14
%% paperName: 2010年海南高考第 14 题 9 分
%% typeId: 4
%% type: 实验题
%% chapter: 功和能
%% point: DIS研究机械能守恒定律
%% method: 无
%% score: 9
%% degree: 450
%% duplicateId: 0
%% body:
利用气垫导轨验证机械能守恒定律，实验装置示意图如图\ref{2010海南14a}所示：

\onepicture[num=1, numstyle=arabic, label=2010海南14a, w=8.5cm]{figs/14a.png}

\begin{enumerate}
	\item
	实验步骤：
	\begin{steps}
		\item
		将气垫导轨放在水平桌面上，桌面高度不低于 $1\Um$，将导轨调至水平；
		\item
		用游标卡尺测量挡光条的宽度 $l$，结果如图\ref{2010海南14b}所示，由此读出 $l=$\tkanswer[1.2cm]{$9.30$}$\Umm$；
		\onepicture[num=2, numstyle=arabic, label=2010海南14b, w=6.5cm, align=right]{figs/14b.png}
		\item
		由导轨标尺读出两光电门中心之间的距离 $s=$\tkanswer[1.5cm]{$60.00$}$\Um$；
		\item
		将滑块移至光电门1左侧某处，待砝码静止不动时，释放滑块，要求砝码落地前挡光条已通过光电门2；
		\item
		从数字计时器（图1中未画出）上分别读出挡光条通过光电门1和光电门2所用的时间 $\Delta t_{1}$ 和 $\Delta t_{2}$；
		\item
		用天平称出滑块和挡光条的总质量 $M$，再称出托盘和砝码的总质量 $m$。
	\end{steps}
	\item
	用表示直接测量量的字母写出下列所示物理量的表达式：
	\begin{steps}
		\item
		滑块通过光电门1和光电门2时瞬时速度分别为 $v_{1}=$\tkanswer[1.8cm]{$\frac{l}{\Delta t_{1}}$}和 $v_{2}=$\tkanswer[1.8cm]{$\frac{l}{\Delta t_{2}}$}。
		\item
		当滑块通过光电门1和光电门2时，系统（包括滑块、挡光条、托盘和砝码）的总动能分别为 $E_{k1}=$\tkanswer[2.8cm]{$\frac{1}{2}(M+m)\frac{l^{2}}{\Delta t_{1}^{2}}$}和 $E_{k2}=$\tkanswer[2.8cm]{$\frac{1}{2}(M+m)\frac{l^{2}}{\Delta t_{2}^{2}}$}。
		\item
		在滑块从光电门1运动到光电门2的过程中，系统势能的减少 $\Delta E_{p}=$\tkanswer[2cm]{$mgs$}（重力加速度为 $g$）。
	\end{steps}
	\item
	如果 $\Delta E_{p}=$\tkanswer[2.5cm]{$E_{k2}-E_{k1}$}，则可认为验证了机械能守恒定律。
\end{enumerate}

%% answer: 
\jdanswer{
\begin{enumerate}
	\item
	② $9.30$，③ $60.00$
	\item
	① $\frac{l}{\Delta t_{1}}$，$\frac{l}{\Delta t_{2}}$；② $\frac{1}{2}(M+m)\frac{l^{2}}{\Delta t_{1}^{2}}$，$\frac{1}{2}(M+m)\frac{l^{2}}{\Delta t_{2}^{2}}$；③ $mgs$
	\item
	$E_{k2}-E_{k1}$
\end{enumerate}
}
%% memo:
\memoanswer{
由于挡光条宽度很小，因此将挡光条通过光电门时的平均速度当作瞬时速度。
}

\item
%% number: 15
%% paperName: 2010年海南高考第 15 题 8 分
%% typeId: 6
%% type: 计算题
%% chapter: 磁场
%% point: 带电粒子在磁场中的运动
%% method: 无
%% score: 8
%% degree: 450
%% duplicateId: 0
%% body:
右图中左边有一对平行金属板，两板相距为 $d$，电压为 $U$；两板之间有匀强磁场，磁感应强度大小为 $B_{0}$，方向与金属板面平行并垂直于纸面朝里。图中右边有一半径为 $R$、圆心为O的圆形区域内也存在匀强磁场，磁感应强度大小为 $B$，方向垂直于纸面朝里。一电荷量为 $q$ 的正离子沿平行于金属板面、垂直于磁场的方向射入平行金属板之间，沿同一方向射出平行金属板之间的区域，并沿直径EF方向射入磁场区域，最后从圆形区域边界上的G点射出。已知弧FG所对应的圆心角为 $\theta$，不计重力。求：

\begin{enumerate}
	\item
	离子速度的大小；
	\item
	离子的质量。
\end{enumerate}

\onepicture[align=right, w=6.5cm]{figs/15.png}

%% answer: 
\jdanswer{
\begin{enumerate}
	\item
	$v=\frac{U}{B_{0}d}$
	\item
	$m=\frac{qBB_{0}Rd}{U}\cot\frac{\theta}{2}$
\end{enumerate}
}
%% memo:
\memoanswer{
（1）由题设知，离子在平行金属板之间做匀速直线运动，它所受到的向上的磁场力和向下的电场力平衡
$qvB_{0}=qE_{0}$，
式中，$v$ 是离子运动速度的大小，$E_{0}$ 是平行金属板之间的匀强电场的强度，有
$E_{0}=\frac{U}{d}$，解得 $v=\frac{U}{B_{0}d}$。

（2）在圆形磁场区域，离子做匀速圆周运动，由洛伦兹力公式和牛顿第二定律有
$qvB=m\frac{v^{2}}{r}$，
式中，$m$ 和 $r$ 分别是离子的质量和它做圆周运动的半径。由题设，离子从磁场边界上的点 G 穿出，离子运动的圆周圆心 $O'$ 必在过 E 点垂直于 EF 的直线上，且在 EG 的垂直平分线上，由几何关系有
$r=R\tan\alpha$，式中，$\alpha$ 是 $OO'$ 与直径 EF 的夹角，由几何关系有
$2\alpha+\theta=\pi$，
解得
$m=\frac{qBB_{0}Rd}{U}\cot\frac{\theta}{2}$。

\onepicture[w=6.0cm]{figs/15_an1.png}

\onepicture[w=6.5cm]{figs/15_an2.png}
}

\item
%% number: 16
%% paperName: 2010年海南高考第 16 题 11 分
%% typeId: 6
%% type: 计算题
%% chapter: 运动和力
%% point: 传送带
%% method: 无
%% score: 11
%% degree: 450
%% duplicateId: 0
%% body:
图\ref{2010海南16a}中，质量为 $m$ 的物块叠放在质量为 $2m$ 的足够长的木板上方右侧，木板放在光滑的水平地面上，物块与木板之间的动摩擦因数为 $\mu=0.2$。在木板上施加一水平向右的拉力 $F$，在 $0\sim3\Us$ 内 $F$ 的变化如图\ref{2010海南16b}所示，图中 $F$ 以 $mg$ 为单位，重力加速度 $g=10\Umsq$。整个系统开始时静止。

\twopicture[labelA=2010海南16a, labelB=2010海南16b, numA=1, numB=2, numstyle=arabic, wA=5.0cm, wB=5.0cm, align=right]
{figs/16a.png}
{figs/16b.png}

\begin{enumerate}
	\item
	求 $1\Us$、$1.5\Us$、$2\Us$、$3\Us$ 末木板的速度以及 $2\Us$、$3\Us$ 末物块的速度；
	\item
	在同一坐标系中画出 $0\sim3\Us$ 内木板和物块的 $v-t$ 图象，据此求 $0\sim3\Us$ 内物块相对于木板滑过的距离。
\end{enumerate}

%% answer: 
\jdanswer{
\begin{enumerate}
	\item
	$v_{1}=4\Ums$，$v_{1.5}=4.5\Ums$，$v_{2}=4\Ums$，$v_{3}=4\Ums$；$v_{2}'=4\Ums$，$v_{3}'=4\Ums$
	\item
	$\Delta s=2.25\Um$
\end{enumerate}
}
%% memo:
\memoanswer{
（1）设木板和物块的加速度分别为 $a$ 和 $a'$，在 $t$ 时刻木板和物块的速度分别为 $v_{t}$ 和 $v_{t}'$，木板和物块之间摩擦力的大小为 $f$，依牛顿第二定律、运动学公式和摩擦定律得
$f=ma'$ ①

$f=\mu mg$（当 $v_{t}'<v_{t}$） ②

$v_{t_{2}}'=v_{t_{1}}'+a'(t_{2}-t_{1})$ ③

$F-f=(2m)a$ ④

$v_{t_{2}}=v_{t_{1}}+a(t_{2}-t_{1})$ ⑤

由①②③④⑤式与题给条件得
$v_{1}=4\Ums$，$v_{1.5}=4.5\Ums$，$v_{2}=4\Ums$，$v_{3}=4\Ums$ ⑥

$v_{2}'=4\Ums$，$v_{3}'=4\Ums$ ⑦

（2）由⑥⑦式得到物块与木板运动的 $v-t$ 图象，如右图所示。在 $0\sim3\Us$ 内物块相对于木板的距离 $\Delta s$ 等于木板和物块 $v-t$ 图线下的面积之差，即图中带阴影的四边形面积，该四边形由两个三角形组成，上面的三角形面积为 $0.25\Um$，下面的三角形面积为 $2\Um$，因此
$\Delta s=2.25\Um$ ⑧

\onepicture[w=6.5cm]{figs/16_an.png}
}

\item
%% number: 17
%% paperName: 2010年海南高考第 17 题 8 分
%% typeId: 6
%% type: 计算题
%% chapter: 气体性质
%% point: 热力学第一定律
%% method: 无
%% score: 8
%% degree: 450
%% duplicateId: 0
%% body:
【模块 3-3 试题】（12 分）

\begin{enumerate}
	\item
	（4 分）下列说法正确的是\tkanswer[1.2cm]{ADE}（填入正确选项前的字母，每选错一个扣2分，最低得分为0分）。

	\fivechoices[answer=ADE]
	{当一定质量的气体吸热时，其内能可能减小}
	{玻璃、石墨和金刚石都是晶体，木炭是非晶体}
	{单晶体有固定的熔点，多晶体和非晶体没有固定的熔点}
	{当液体与大气相接触时，液体表面层内的分子所受其它分子作用力的合力总是指向液体内部}
	{气体分子单位时间内与单位面积器壁碰撞的次数，与单位体积内气体的分子数和气体温度有关}
	\item
	（8 分）如右图，体积为 $V$、内壁光滑的圆柱形导热气缸顶部有一质量和厚度均可忽略的活塞；气缸内密封有温度为 $2.4T_{0}$、压强为 $1.2p_{0}$ 的理想气体。$p_{0}$ 和 $T_{0}$ 分别为大气的压强和温度。已知：气体内能 $U$ 与温度 $T$ 的关系为 $U=\alpha T$，$\alpha$ 为正的常量；容器内气体的所有变化过程都是缓慢的。求：

	\begin{enumerate}[label = (\roman*)]
		\item
		气缸内气体与大气达到平衡时的体积 $V_{1}$；
		\item
		在活塞下降过程中，气缸内气体放出的热量 $Q$。
	\end{enumerate}

	\onepicture[align=right, w=2.8cm]{figs/17.png}
\end{enumerate}

%% answer: 
\jdanswer{
\begin{enumerate}
	\item
	ADE
	\item
	（i）$V_{1}=\frac{1}{2}V$；（ii）$Q=\frac{1}{2}p_{0}V+\alpha T_{0}$
\end{enumerate}
}
%% memo:
\memoanswer{
（1）一定质量的气体吸热时，如果同时对外做功，且做的功大于吸收的热量，则内能减小，（A）正确；玻璃是非晶体，（B）错；多晶体也有固定的熔点，（C）错；液体表面层内的分子与液体内部分子间距离、密度都大于大气，因此分子力的合力指向液体内部，（D）正确；气体分子单位时间内与单位面积器壁碰撞的次数，决定气体的压强，因此与单位体积内分子数和气体的温度有关，（E）对。

（2）（i）在气体由压缩 $p=1.2p_{0}$ 下降到 $p_{0}$ 的过程中，气体体积不变，温度由 $T=2.4T_{0}$ 变为 $T_{1}$，由查理定律得 $\frac{T_{1}}{T}=\frac{p_{0}}{p}$ ①

在气体温度由 $T_{1}$ 变为 $T_{0}$ 的过程中，体积由 $V$ 减小到 $V_{1}$，气体压强不变，由盖·吕萨克定律得
$\frac{V}{V_{1}}=\frac{T_{1}}{T_{0}}$ ②

由①②式得 $V_{1}=\frac{1}{2}V$ ③

（ii）在活塞下降过程中，活塞对气体做的功为
$W=p_{0}(V-V_{1})$ ④

在这一过程中，气体内能的减少为
$\Delta U=\alpha(T_{1}-T_{0})$ ⑤

由热力学第一定律得，气缸内气体放出的热量为
$Q=W+\Delta U$ ⑥

由②③④⑤⑥式得
$Q=\frac{1}{2}p_{0}V+\alpha T_{0}$ ⑦
}

\item
%% number: 18
%% paperName: 2010年海南高考第 18 题 6 分
%% typeId: 6
%% type: 计算题
%% chapter: 机械波
%% point: 干涉衍射
%% method: 无
%% score: 6
%% degree: 450
%% duplicateId: 0
%% body:
【模块 3-4 试题】（12 分）

\begin{enumerate}
	\item
	（6 分）一光线以很小的入射角 $i$ 射入一厚度为 $d$、折射率为 $n$ 的平板玻璃，求出射光线与入射光线之间的距离（$i$ 很小时，$\sin i=i$，$\cos i=1$）。
	\item
	（6 分）右图为某一报告厅主席台的平面图，AB是讲台，$S_{1}$、$S_{2}$是与讲台上话筒等高的喇叭，它们之间的相互位置和尺寸如图所示。报告者的声音放大后经喇叭传回话筒再次放大时可能会产生啸叫。为了避免啸叫，话筒最好摆放在讲台上适当的位置，在这些位置上两个喇叭传来的声音因干涉而相消。已知空气中声速为 $340\Ums$，若报告人声音的频率为 $136\UHz$，问讲台上这样的位置有多少个？

	\onepicture[align=right, w=6.5cm]{figs/18.png}
\end{enumerate}

%% answer: 
\jdanswer{
\begin{enumerate}
	\item
	$\frac{(n-1)id}{n}$
	\item
	4
\end{enumerate}
}
%% memo:
\memoanswer{
（1）如图，设光线以很小的入射角 $i$ 入射到平板玻璃表面上的 A 点，折射角为 $\gamma$，从平板玻璃另一表面上的 B 点射出。设 AC 为入射光线的延长线。由折射定律和几何关系可知，它与出射光线平行。过 B 点作 $BD\perp AC$，交 AC 于 D 点，则 BD 的长度就是出射光线与入射光线之间的距离，由折射定律得
$\frac{\sin i}{\sin\gamma}=n$ ①

由几何关系得 $\angle BAD=i-\gamma$ ②

$\overline{AB}=\frac{d}{\cos\gamma}$ ③

出射光线与入射光线之间的距离为
$\overline{BD}=\overline{AB}\sin(i-\gamma)$ ④

当入射角 $i$ 很小时，有
$\sin i=i$，$\sin\gamma=\gamma$，$\sin(i-\gamma)=i-\gamma$，$\cos\gamma=1$

由此及①②③④式得 $\overline{BD}=\frac{(n-1)d}{n}i$ ⑤

\onepicture[w=3.2cm]{figs/18_an1.jpg}

（2）相应于声频 $f=136\UHz$ 的声波的波长是
$\lambda=\frac{v}{f}=2.5\Um$ ①

式中 $v=340\Ums$ 是空气中的声速。在右图中，O 是 AB 的中点，P 是 OB 上任一点。将 $S_{1}P-S_{2}P$ 表示为
$S_{1}P-S_{2}P=k\frac{\lambda}{2}$ ②

式中 $k$ 为实数，当 $k=0,2,4\cdots$ 时，从两个喇叭来的声波因干涉而加强；当 $k=1,3,5\cdots$ 时，从两个喇叭来的声波因干涉而相消。由此可知，O 是干涉加强点；对于 B 点，
$S_{1}B-S_{2}B=(20-15)\Um=4\frac{1}{2}\lambda$ ③

所以，B 点也是干涉加强点。因而 O、B 之间有两个干涉相消点，由对称性可知，AB 上有 4 个干涉相消点。

\onepicture[w=7.0cm]{figs/18_an2.jpg}
}

\item
%% number: 19
%% paperName: 2010年海南高考第 19 题 8 分
%% typeId: 6
%% type: 计算题
%% chapter: 运动和力
%% point: 动量和能量
%% method: 无
%% score: 8
%% degree: 450
%% duplicateId: 0
%% body:
【模块 3-5 试题】（12 分）

\begin{enumerate}
	\item
	（4 分）能量为 $E_{i}$ 的光子照射基态氢原子，刚好可使该原子中的电子成为自由电子，这一能量 $E_{i}$ 称为氢的电离能。现用一频率为 $\nu$ 的光子从基态氢原子中击出了一电子，该电子在远离核后速度大小为\tkanswer[3cm]{$\sqrt{\frac{2(h\nu-E_{i})}{m}}$}（用光子频率 $\nu$、电子质量 $m$、氢的电离能 $E_{i}$ 与普朗克常量 $h$ 表示）
	\item
	（8 分）在核反应堆中，常用减速剂使快中子减速。假设减速剂的原子核质量是中子的 $k$ 倍，中子与原子核的每次碰撞都可看成是弹性正碰。设每次碰撞前原子核可认为上静止的，求 $N$ 次碰撞后中子速率与原速率之比。
\end{enumerate}

%% answer: 
\jdanswer{
\begin{enumerate}
	\item
	$\sqrt{\frac{2(h\nu-E_{i})}{m}}$
	\item
	$\left(\frac{k-1}{k+1}\right)^{N}$
\end{enumerate}
}
%% memo:
\memoanswer{
（1）由能量守恒得 $\frac{1}{2}mv^{2}=h\nu-E_{i}$，解得电子速度为 $v=\sqrt{\frac{2(h\nu-E_{i})}{m}}$。

（2）设中子和作减速剂的物质的原子核 A 的质量分别为 $m_{\mathrm{n}}$ 和 $m_{A}$，碰撞后速度分别为 $v_{\mathrm{n}}'$ 和 $v_{A}'$，碰撞前后的总动量和总能量守恒，有
$m_{\mathrm{n}}v_{\mathrm{n}}=m_{\mathrm{n}}v_{\mathrm{n}}'+m_{A}v_{A}'$ ①

$\frac{1}{2}m_{\mathrm{n}}v_{\mathrm{n}}^{2}=\frac{1}{2}m_{\mathrm{n}}v_{\mathrm{n}}'^{2}+\frac{1}{2}m_{A}v_{A}'^{2}$ ②

式中 $v_{\mathrm{n}}$ 为碰撞前中子速度，由题设
$m_{A}=km_{\mathrm{n}}$ ③

由①②③式得，经 1 次碰撞后中子速率与原速率之比为
$\frac{\left|v_{\mathrm{n}}'\right|}{\left|v_{\mathrm{n}}\right|}=\frac{k-1}{k+1}$ ④

经 N 次碰撞后，中子速率与原速率之比为
$\left(\frac{k-1}{k+1}\right)^{N}$ ⑤
}

\end{enumerate}

\end{document}
